\section{State Estimation and Observers}
\ptitle{State Estimation and Observer Feedback}

Estimate the state $\hat{x}(t)$ and apply $u(t) = u(\hat{x}(t))$.

\subsection{Introduction}
\subsubsection{LTI Observers}
\paragraph{Noise-Free Case}
\ptitle{System}
\begin{align*}
    \dot{x} & = Ax + Bu \\
    y       & = Cx
\end{align*}

\newpar{}
\ptitle{Observer}
\begin{align*}
    \dot{\hat{x}} & = A\hat{x} + Bu + L(y - C\hat{x})                  \\
                  & = (A - LC)\hat{x} + Bu + Ly \quad \text{(useful?)}
\end{align*}

\newpar{}
\ptitle{Error Dynamics}

Define the error $e = x - \hat{x}$
\begin{align*}
    \dot{e} & = (Ax + Bu) - (A - LC)\hat{x} - Bu - LCx \\
            & = (A - LC)e
\end{align*}

\newpar{}
\ptitle{Closed-Loop Stability}

Use control $u = -K\hat{x} = -K(x - e)$:
\begin{equation*}
    \begin{bmatrix}
        \dot{x} \\
        \dot{e}
    \end{bmatrix}
    =
    \begin{bmatrix}
        A - BK & BK     \\
        0      & A - LC
    \end{bmatrix}
    \begin{bmatrix}
        x \\
        e
    \end{bmatrix}
\end{equation*}

If $A - BK$ and $A - LC$ are stable, then the closed-loop is stable. Hence, by the \textbf{separation principle} we can design observer and controller \textbf{separately} to obtain a stable CL.

\newpar{}
\ptitle{Limitations}

\begin{itemize}
    \item The separation principle does not hold in general for nonlinear systems.
    \item For the nonlinear case, the following methods are common
          \begin{itemize}
              \item Extended Kalman Filter (EKF)
              \item Recursive Prediction Error (RPE)
              \item Unscented Kalman Filter (UKF)
              \item Moving Horizon Estimation (MHE)
              \item High Gain Observers (HGO)
          \end{itemize}
\end{itemize}
\paragraph{Noisy Case}

\ptitle{System}
\begin{align*}
    x(k{+}1) & = Ax(k) + Bu(k) + L_w w(k) \\
    y(k)     & = Cx(k) + v(k)
\end{align*}

$w(k) \sim \mathcal{N}(0, Q(k))$, \quad $v(k) \sim \mathcal{N}(0, R(k))$

\newpar{}
\ptitle{Observer}
\begin{align*}
    \hat{x}(k{+}1) & = A\hat{x}(k) + Bu(k) + L\big(y(k) - C\hat{x}(k) \big) \\
    \hat{y}(k)     & = C\hat{x}(k)
\end{align*}

\newpar{}
\ptitle{Error Dynamics}

Define the error $e(k) = x(k) - \hat{x}(k)$:
\begin{align*}
    e(k{+}1) & = (A - L C)e(k) + \underbrace{L_w w(k) - L v(k)}_{:=\rho(k)\text{ (Gaussian noise)}} \\
\end{align*}

\newpar{}
\ptitle{Closed-Loop Stability}

With control $u(k) = -K\hat{x}(k) = -K(x(k) - e(k))$, we have:
\begin{align*}
    x(k{+}1) & = (A - BK)x(k) + BKe(k) + L_w w(k)
\end{align*}

The combined system becomes:
{\small
\begin{equation*}
    \begin{bmatrix}
        x(k{+}1) \\
        e(k{+}1)
    \end{bmatrix}
    =
    \begin{bmatrix}
        A - BK & BK      \\
        0      & A - L C
    \end{bmatrix}
    \begin{bmatrix}
        x(k) \\
        e(k)
    \end{bmatrix}
    +
    \begin{bmatrix}     % TODO: Doesn't this matrix lack an entry?
        L_w & 0  \\
        0   & -L
    \end{bmatrix}
    \begin{bmatrix}
        w(k) \\
        v(k)
    \end{bmatrix}
\end{equation*}
}

\newpar{}
\ptitle{Remarks}
\begin{itemize}
    \item The error dynamics are now driven by Gaussian noise.
    \item If $A - BK$ and $A - L C$ are stable, the system rejects noise over time.
    \item The separation principle still applies in this linear stochastic setting.
\end{itemize}


\subsubsection{Estimation of Constant Vectors}
In this section we look at the estimation of a \textit{constant} vector from \textit{repeated observations}.
\newpar{}
\ptitle{System}
\begin{equation*}
    y = Cx + w, \quad x \in \mathbb{R}^n, \quad y, w \in \mathbb{R}^p
\end{equation*}
with $w$ being white noise.
\newpar{}
Our goal is to compute $\hat{x}$ from given measurements.

\newpar{}
\ptitle{Definitions}
\begin{itemize}
    \item measurement estimate: $\hat{y} = C \hat{x}$
    \item residual: $\varepsilon_y = \hat{y} - y$
    \item cost function:
          \begin{equation*}
              J[\hat{x}] = \frac{1}{2} \varepsilon_y^\top \varepsilon_y
          \end{equation*}
\end{itemize}

\newpar{}
\ptitle{Optimal Estimate}

As the cost is a simple quadratic form, the optimal estimate is given by pseudo inversion
\begin{equation*}
    \hat{x} = {(C^\top C)}^{-1} C^\top y
\end{equation*}
\newpar{}
\ptitle{Practical Considerations}

This static approach introduces averaging. Hence, one must ensure that no undesired averaging happens, e.g.\ of transient and non-transient phases.

\paragraph{Least Squares for Multiple Measurements}
Consider a process modelled by the simple linear relationship
\begin{equation*}
    y = Cx+w
\end{equation*}
The actual value of $x$ needs to be reconstructed from the data
\begin{equation*}
    Y = [y_1;y_2; \dots ;y_N] = [Cx_1+w_1; Cx_2+w_2; \dots ;Cx + w_N]
\end{equation*}
The optimal estimate $\hat{x}$ is given by
\begin{equation*}
    \hat{x} = H^{\dagger}Y, \quad H = [C;C; \dots ;C]
\end{equation*}
where $H^{\dagger} = {(H^\top H)}^{-1}H^\top$.

The estimate $\hat{x}$ approaches the true value as $N$ grows.

\paragraph{Weighted Least Squares}

Assign weight matrix $N^{-1}$ to residuals:
\begin{equation*}
    \varepsilon_y = N^{-1}(\hat{y} - y)
\end{equation*}

Then minimize:
\begin{equation*}
    \hat{x} = {(C^\top S^{-1} C)}^{-1} C^\top S^{-1} y, \quad S = N^{-\top} N^{-1}
\end{equation*}

This allows tuning the significance of different measurements. Example: Give less weight to older measurements.

\subsection{Kalman Filter}
In this section we look at the estimation of dynamical vectors.
\ptitle{Kalman Filter Problem Setup}
\begin{align*}
    x(k + 1) & = Ax(k) + Bu(k) + Lw(k), \quad x(0) = x_0 \\
    y(k)     & = Cx(k) + v(k)
\end{align*}
With assumptions on the random variables
\begin{align*}
     & E[x_0] = \bar{x}_0, & E[(x_0 - \bar{x}_0){(x_0 - \bar{x}_0)}^\top] = P(0) \\
     & E[w(k)] = 0,        & E[w(k){w(k)}^\top] = Q(k)                           \\
     & E[v(k)] = 0,        & E[v(k){v(k)}^\top] = R(k)
\end{align*}
We also have the following assumptions on uncorrelatedness
\begin{gather*}
    E[v(k){v(j)}^\top] = 0,\;E[w(k){w(j)}^\top] = 0,\quad k\neq j \\
    E[v(k){w(j)}^\top] = 0,\quad \forall k, j
\end{gather*}

\subsubsection{Algorithm}

\begin{enumerate}
    \item \textbf{Prior Update (Prediction)}
          \begin{align*}
              P(k + 1|k) & = A P(k|k) A^\top + L Q(k) L^\top \\
              x(k + 1|k) & = A x(k|k) + Bu(k)
          \end{align*}
    \item \textbf{Posterior Update (Measurement Update)}
          \begin{equation*}
              K(k + 1) = P(k + 1|k) C^\top {\left(C P(k + 1|k) C^\top + R(k)\right)}^{-1}
          \end{equation*}
          yielding
          \begin{equation*}
              P(k + 1|k + 1) = (I - K(k + 1) C) P(k + 1|k)
          \end{equation*}
          {\small
          \begin{equation*}
              x(k + 1|k + 1) = x(k + 1|k) + K(k + 1) \left(y(k + 1) - C x(k + 1|k)\right)
          \end{equation*}
          }
\end{enumerate}

\textbf{Note} that $K$ is chosen such that the trace of $P(k+1)$ is minimized, yielding minimal uncertainty in the state estimation error.

\subsection{Extended Kalman Filter (EKF) and EKF with Parameter Estimation}
\subsubsection{EKF}
\ptitle{System}

The Extended Kalman Filter (EKF) adapts the standard Kalman Filter to nonlinear systems of the form:
\begin{align*}
    x(k+1)                               & = f(x(k), u(k)) + w(k)         \\
    y(t_k)                               & = h(x(t_k)) + v(t_k)           \\
    w(k) \sim \mathcal{N}(0, Q(k)),\quad & v(k) \sim \mathcal{N}(0, R(k))
\end{align*}

\newpar{}
\ptitle{Key Ideas}

\begin{enumerate}
    \item Propagate the state and covariance between measurements using the nonlinear model
    \item Linearize the system at the current estimate $\hat{x}(t)$ to compute updates at discrete times $t_k$
\end{enumerate}

\paragraph{Algorithm}

\newpar{}
\ptitle{Prior Update (Prediction)}

Between $t_k$ and $t_{k+1}$: Given $\hat{x}(t_k^+), P(t_k^+)$, propagate forward using the nonlinear model:
\begin{align*}
    \hat{x}(t_{k+1}^-) & = f(\hat{x}(t_k^+), u(t_k))            \\
    P(t_{k+1}^-)       & = F(t_k) P(t_k^+) F(t_k)^\top + Q(t_k)
\end{align*}
with the Jacobian:
\begin{equation*}
    F(t_k) = \frac{\partial f}{\partial x} \bigg|_{x = \hat{x}(t_k^+)}
\end{equation*}

\newpar{}
\ptitle{Posterior Update (Measurement Update)}

At $t = t_{k+1}$: Define the innovation
\begin{equation*}
    e(t_{k+1}) = y(t_{k+1}) - h(\hat{x}(t_{k+1}^-)) \text{ (innovation)}
\end{equation*}
and obtain the linearization
\begin{align*}
    H(t_{k+1}^-) & = \frac{\partial h}{\partial x} \bigg|_{x = \hat{x}(t_{k+1}^-)} \\
    A(t_{k+1})   & = R + H(t_{k+1}^-) P(t_{k+1}^-) {H(t_{k+1}^-)}^\top             \\
    K(t_{k+1})   & = P(t_{k+1}^-) {H(t_{k+1}^-)}^\top {A(t_{k+1})}^{-1}
\end{align*}
This yields the measurement update
\begin{align*}
    \hat{x}(t_{k+1}^+) & = \hat{x}(t_{k+1}^-) + K(t_{k+1}) e(t_{k+1})             \\
    P(t_{k+1}^+)       & = P(t_{k+1}^-) - K(t_{k+1}) A(t_{k+1}) {K(t_{k+1})}^\top
\end{align*}

\subsubsection{Hybrid EKF}
In practice, the system evolves continuously, but measurements are taken at discrete times $t_k$.
\begin{align*}
    \dot{x}(t) & = f(x(t)) + w(t)     \\
    y(t_k)     & = h(x(t_k)) + v(t_k)
\end{align*}

\paragraph{Algorithm}

\ptitle{Prior Update (Prediction)}

For $t \in (t_k, t_{k+1})$: Given $\hat{x}(t_k^+)$, propagate between measurements
\begin{align*}
    \dot{\hat{x}}(t) & = f(\hat{x}(t))                    \\
    \dot{P}(t)       & = F(t) P(t) + P(t) {F(t)}^\top + Q
\end{align*}
with Jacobian:
\begin{equation*}
    F(t) = \frac{\partial f}{\partial x} \bigg|_{x = \hat{x}(t)}
\end{equation*}
which yields $\hat{x}(t_{k+1}^-), P(t_{k+1}^-)$ for the next step.

\newpar{}
\ptitle{Posterior Update (Measurement Update)}

The posterior update is the same as in the standard EKF.

\subsubsection{KF with Parameter Estimation}
\ptitle{System}

To estimate unknown parameters $p$, augment the state:
\begin{equation*}
    z(t) =
    \begin{bmatrix}
        x(t) \\
        p(t)
    \end{bmatrix}
\end{equation*}
Augmented system with $p$ constant but uncertain:
\begin{align*}
    \dot{x}(t) & = f(x(t), u(t), p(t)) + w_1(t) \\
    \dot{p}(t) & = 0 + w_2(t)                   \\
    y(t_k)     & = h(x(t_k), p(t_k)) + v(t_k)
\end{align*}
with covariances
\begin{align*}
    Q & = \mathbb{E}[w_1 w_1^\top], \quad Q_{pp} = \mathbb{E}[w_2 w_2^\top], \quad R = \mathbb{E}[v v^\top]
\end{align*}

\newpar{}
\ptitle{Prior Update (Prediction)}

Covariance propagation:
\begin{align*}
    \dot{P}_{xx} & = A P_{xx} + P_{xx} A^\top + G P_{xp}^\top + P_{xp} G^\top + Q \\
    \dot{P}_{xp} & = A P_{xp} + G P_{pp}                                          \\
    \dot{P}_{pp} & = Q_{pp}
\end{align*}

\newpar{}
\ptitle{Posterior Update (Measurement Update)}
Innovation:
\begin{gather*}
    e(t_{k+1}) = y(t_{k+1}) - \hat{y}(t_{k+1})\\
    \hat{y}(t_{k+1}) = h(\hat{x}(t_{k+1}^-), u(t_{k+1}), \hat{p}(t_{k+1}^-))
\end{gather*}
\newpar{}
Linearization:
\begin{align*}
    K(t_{k+1}) & = \left( P_{xx}(t_{k+1}^-) \frac{\partial h}{\partial x}^\top + P_{xp}(t_{k+1}^-) \frac{\partial h}{\partial p}^\top \right) A(t_{k+1})^{-1}      \\
    L(t_{k+1}) & = \left( P_{xp}(t_{k+1}^-)^\top \frac{\partial h}{\partial x}^\top + P_{pp}(t_{k+1}^-) \frac{\partial h}{\partial p}^\top \right) A(t_{k+1})^{-1}
\end{align*}
Posterior update:
\begin{align*}
    \hat{x}(t_{k+1}^+) & = \hat{x}(t_{k+1}^-) + K(t_{k+1}) e(t_{k+1}) \\
    \hat{p}(t_{k+1})   & = \hat{p}(t_k) + L(t_{k+1}) e(t_{k+1})
\end{align*}


\subsection{Moving Horizon Estimation (MHE)}
\begin{center}
    \includegraphics[width=0.8\linewidth]{12_MHE.png}
\end{center}

Moving Horizon Estimation (MHE) is based on optimization over a moving time window to estimate system states.

\begin{itemize}
    \item Collect recent inputs $u_M(\tau)$ and outputs $y_M(\tau)$ over the interval $[t - M, t]$
    \item Minimize a cost function involving model prediction errors and measurement residuals
\end{itemize}
The idea is to calculate the trajectory of states $x(\tau), \tau \in [t-M : t]$ that better ``explains'' the given history and the model.

\subsubsection{Standard MHE Formulation}

Given a nonlinear system:
\begin{align*}
    x_{k+1} & = f(x_k, u_k) + w_k \\
    y_k     & = h(x_k) + v_k
\end{align*}

With measurements $y_{t-M}, \dots, y_t$ and known inputs $u_{t-M}, \dots, u_{t-1}$, the MHE estimates the trajectory by solving the optimization problem:
\begin{gather*}
    \min_{\{{x_k}\}_{k=t-M}^{t}} \|x_{t-M} - \bar{x}_{t-M}\|^2_{P^{-1}}\\ % ChkTex 3
    + \sum_{k=t-M}^{t-1} \|x_{k+1} - f(x_k, u_k)\|^2_{Q^{-1}} + \sum_{k=t-M}^{t} \|y_k - h(x_k)\|^2_{R^{-1}}
\end{gather*}

Subject to:
\begin{itemize}
    \item $x_k \in \mathcal{X}$: state constraints (optional)
    \item $\bar{x}_{t-M}$: prior state (arrival cost) with covariance $P$
    \item $Q$, $R$: process and measurement noise covariances
\end{itemize}
The tuning parameters characterize the degree uf trust into the mesurements and the model.

\paragraph{Properties and Benefits of MHE}

\begin{itemize}
    \item [+] Can explicitly enforce constraints (e.g.\ energy balances, physical bounds). Reduces ambiguity and improves estimation quality.
    \item [+] Well suited to nonlinear systems or when constraints are critical
    \item [+] Robust
    \item Converges to the Kalman Filter as $M \to \infty$ in linear systems
\end{itemize}

\newpar{}
\ptitle{Remarks on MHE and Stability}

Stability of MHE and MPC requires:
\begin{enumerate}
    \item MHE must converge rapidly (i.e.\ tight error bounds after few iterations)
    \item MPC must maintain a continuous value function (robustness to perturbations)
\end{enumerate}

\subsection{High Gain Observers (Not in Lecture!)}

\ptitle{System}

Consider the system of the \textit{form}:
\begin{align*}
    \dot{x}_1 & = x_2        \\
    \dot{x}_2 & = \phi(x, u) \\
    y         & = x_1
\end{align*}

Assume a locally Lipschitz state feedback $u = \gamma(x)$ stabilizes the origin $x = 0$ of the closed-loop system:
\begin{align*}
    \dot{x}_1 & = x_2                \\
    \dot{x}_2 & = \phi(x, \gamma(x)) \\
    y         & = x_1
\end{align*}
\newpar{}
\ptitle{Observer}

To implement this control using only output $y$, define the observer:
\begin{align*}
    \dot{\hat{x}}_1 & = \hat{x}_2 + h_1(y - \hat{x}_1)                        \\
    \dot{\hat{x}}_2 & = \phi_0(\hat{x}, \gamma(\hat{x})) + h_2(y - \hat{x}_1)
\end{align*}

where $\phi_0(x, u)$ is a nominal model of $\phi(x, u)$.

\newpar{}
\ptitle{Estimation Error}

Define the estimation error $\tilde{x} = x - \hat{x}$, then:
\begin{align*}
    \dot{\tilde{x}}_1 & = -h_1 \tilde{x}_1 + \tilde{x}_2        \\
    \dot{\tilde{x}}_2 & = -h_2 \tilde{x}_1 + \delta(x, \hat{x})
\end{align*}
with
\begin{equation*}
    \delta(x, \hat{x}) = \phi(x, \gamma(x)) - \phi_0(\hat{x}, \gamma(\hat{x}))
\end{equation*}

\newpar{}
\ptitle{Stable Estimation Error}

We aim to choose $(h_1, h_2)$ such that $\tilde{x}(t) \to 0$. In the ideal case $\delta \equiv 0$, this can be done by making the matrix
\begin{equation*}
    A_0 =
    \begin{bmatrix}
        -h_1 & 1 \\
        -h_2 & 0
    \end{bmatrix}
\end{equation*}
stable.

\newpar{}
If $\delta \neq 0$, we minimize the gain of the transfer function $G_0 : \delta \mapsto \tilde{x}$. It takes the form:
\begin{equation*}
    G_0 =
    \frac{1}{s^2 + h_1 s + h_2}
    \begin{bmatrix}
        1 \\
        s + h_1
    \end{bmatrix}
\end{equation*}
This gain can be made arbitrarily small by setting $h_2 \gg h_1 \gg 1$. For instance, with
\begin{equation*}
    h_1 = \frac{\alpha_1}{\varepsilon}, \quad h_2 = \frac{\alpha_2}{\varepsilon^2}, \quad \varepsilon \ll 1
\end{equation*}
we have
\begin{equation*}
    \lim_{\varepsilon \to 0} \|G_0(s)\| = 0
\end{equation*}

\newpar{}
\ptitle{Peaking Phenomenon}

However, the $(2,1)$ entry of $\exp(A_0 t)$ grows with $h_1, h_2$, causing a peaking phenomenon. That is, $\tilde{x}(t)$ becomes arbitrarily large temporarily, which may result in a large control input $u = \gamma(\hat{x})$, potentially pushing the system outside the region of attraction (unstable error dynamics).
\newpar{}
Solution: \textbf{saturating} the control action. Since the peaking is short-lived, we assume that $\hat{x}(t)$ converges faster than the system diverges, and the resulting feedback remains stabilizing.

\subsubsection{SISO HGO Theorem}
\ptitle{System}
\begin{align*}
    \dot{x} & = A_C x + B_C \phi(x, u) \\
    y       & = C_C x
\end{align*}
\newpar{}
\ptitle{Controller}

Assume there exists a stabilizing state feedback controller:
\begin{align*}
    \dot{v} & = \Gamma(v, x) \\
    u       & = \gamma(v, x)
\end{align*}
with $\Gamma$ and $\gamma$ globally bounded.

Now consider the output feedback implementation:
\begin{align*}
    \dot{v}       & = \Gamma(v, \hat{x})                                        \\
    \dot{\hat{x}} & = A_C \hat{x} + B_C \phi_0(\hat{x}, u) + H(y - C_C \hat{x}) \\
    u             & = \gamma(v, \hat{x})
\end{align*}

with ($\rho$ is the relative degree)
\begin{equation*}
    H =
    \begin{bmatrix}
        \frac{\alpha_1}{\varepsilon}   \\
        \frac{\alpha_2}{\varepsilon^2} \\
        \vdots                         \\
        \frac{\alpha_\rho}{\varepsilon^\rho}
    \end{bmatrix}
    \in \mathbb{R}^\rho
\end{equation*}
and $\varepsilon>0$ and choose $\alpha_i$ such that
\begin{equation*}
    s^\rho + \alpha_1 s^{\rho - 1} + \cdots + \alpha_\rho = 0
\end{equation*}
has all its roots in the complex left half plane.
\newpar{}
Then, there exists $\varepsilon_0 > 0$ such that for all $\varepsilon < \varepsilon_0$, the closed-loop system is \textbf{exponentially stable}.

\newpar{}
\textbf{Remark}

This formulation neglects unobserved zero dynamics. If these exist, we assume the original system is minimum phase. Then the stability result holds.


